\section{Related Work}

\textbf{Robustness and Bias in Vision Models:}
Although modern deep learning models encode overrepresented training patterns as implicit priors, these task-extrinsic biases often remain hidden under standard benchmarks until those assumptions are violated.
Extensive work has demonstrated this across visual recognition, including biases toward texture over shape, shortcut correlations, photographer preferences, object position, and background context~\citep{geirhos2018imagenet, geirhos2020shortcut, torralba2011unbiased, xiao2020noise, liu2026optigeoefficientmonoculargeometry}.
These findings have motivated robustness benchmarks beyond clean in-distribution accuracy~\citep{hendrycks2019benchmarking} and approaches based on architectural inductive biases or inference-time adaptation~\citep{cohen2016group, weiler2019general, sun2020test, wang2020tent}.
Despite this extensive literature, robustness has been studied primarily in recognition tasks such as classification, segmentation, and detection. 
In single-image geometry, prior studies have identified specific camera pose and spatial distribution biases in depth estimation ~\citep{vandijk2019neuralnetworksdepthsingle, zhao2021cameraposemattersimproving}, while recent work highlights broader viewpoint instabilities in vision foundation models~\citep{michalkiewicz2024viewscreatedequalanalyzing}. Similarly, \cite{jiang2025youviewpointinvariant} investigates viewpoint invariance in policy learning.
However, whether similar priors arise in feed-forward geometric foundation models remains largely unexplored.

\textbf{Image Orientation and Uprightness Estimation:}
The resulting sensitivity to image orientation naturally raises the question of whether canonical upright orientation can be estimated reliably at inference time.
Existing approaches include geometric methods based on line classification~\citep{ctrlc, lee2021ctrl}, feature matching \citep{nordstrom2026whohandlesorientationcvpr, nordstrom2026lomawhohandlesorientationeccv} and learning-based methods such as UprightNet~\citep{xian2019uprightnet}, which predicts 2-DoF camera orientation from indoor surface normals, and GeoCalib~\citep{veicht2024geocalib}, which jointly estimates camera intrinsics and gravity direction through differentiable geometric optimization.
Inference-time canonicalization offers an alternative for robustness under distribution shifts. 
FOCAL~\citep{singhal2025testtimeimageCanonicalization} searches for a canonical input transformation using test time optimization of likelihoods with foundation-model, while G3TUP!~\citep{kani2026g3t} fine-tunes VGGT to produce upright point clouds. 

In contrast, we identify upright bias as a previously unrecognized failure mode shared across modern feed-forward reconstruction models and introduce ROVER, a Plug-and-Play inference-time canonicalization framework that estimates orientation using DINOv2, requiring neither scene-specific geometric structure nor model retraining and generalizing across indoor and outdoor scenes.

% {\color{blue}
% \begin{enumerate}

%     \item \textbf{How Do Neural Networks See Depth in Single Images?}~\citep{vandijk2019neuralnetworksdepthsingle} --- 
%     \textit{Tom van Dijk, Guido C. H. E. de Croon} --- \textbf{ICCV 2019}.  
%     The paper finds that MonoDepth relies strongly on an object's \textbf{vertical image position} rather than apparent size for depth estimation. Because this cue depends on camera pose, changes in pitch can lead to systematic depth errors.

%     \item \textbf{Camera Pose Matters: Improving Depth Prediction by Mitigating Pose Distribution Bias}~\citep{zhao2021cameraposemattersimproving} --- 
%     \textit{Yunhan Zhao, Shu Kong, Charless Fowlkes} --- \textbf{CVPR 2021}.  
%     The authors identify a \textbf{camera-pose distribution bias} in monocular depth estimation, where models become unreliable on uncommon or unseen viewpoints. They mitigate this using perspective-aware augmentation and explicit camera-pose conditioning.

%     % \item \textbf{Breaking the Horizontal Prior: From Long-Tailed Orientation Bias to Roll-Robust Monocular Depth Estimation}~\citep{tang2026breakinghorizontalpriorlongtailed} --- 
%     % \textit{Kaihua Tang et al.} --- \textbf{2026}.  
%     % This work identifies a \textbf{Horizontal Prior} caused by the dominance of upright images in monocular depth training data, leading to degradation under camera roll. The authors propose an Invariant Depth Constraint to learn more rotation-stable representations.

%     \item \textbf{AerialMegaDepth: Learning Aerial-Ground Reconstruction and View Synthesis}~\citep{vuong2025aerialmegadepthlearningaerialgroundreconstruction} --- 
%     \textit{Khiem Vuong, Anurag Ghosh, Deva Ramanan, Srinivasa Narasimhan, Shubham Tulsiani} --- \textbf{CVPR 2025}.  
%     The paper shows that learned 3D reconstruction struggles under \textbf{extreme aerial-ground viewpoint changes}, partly due to insufficient viewpoint diversity in training data. Fine-tuning on a dedicated aerial-ground corpus substantially improves camera estimation and registration.


%     \item \textbf{Not All Views are Created Equal: Analyzing Viewpoint Instabilities in Vision Foundation Models}~\citep{michalkiewicz2024viewscreatedequalanalyzing} --- 
%     \textit{Mateusz Michalkiewicz et al.} --- \textbf{ICCV 2025}.  
%     The paper analyzes \textbf{viewpoint instability} in foundation models and shows that representation quality can vary significantly with viewing angle. These instabilities also affect downstream 3D reasoning and reconstruction tasks.

%     \item \textbf{OptiGeo: Efficient Monocular Geometry for Embodied Perception in Optically Challenging Scenes}~\citep{liu2026optigeoefficientmonoculargeometry} --- 
%     \textit{Muxin Liu et al.} --- \textbf{2026}.  
%     OptiGeo shows that failures on transparent and reflective surfaces can arise from \textbf{biased real-depth supervision}, where sensor errors are inherited by learned geometry models. It mitigates this by replacing unreliable supervision with clean-geometry teacher and residual-trimmed alignment.

% \end{enumerate}
% }